O problema que se nos apresenta hoje é uma joia da Álgebra Linear, uma proposição que tece uma
conexão profunda e elegante entre três conceitos centrais: o espectro de autovalores de um
operador, sua estrutura geométrica (diagonalizabilidade) e sua dinâmica (a existência de um vetor
cíclico).

Vamos provar esta equivalência com o rigor e a clareza que caracterizam nossa disciplina,
demonstrando as duas implicações da proposição "se e somente se".

\textbf{Enunciado do Problema:} \\
Seja $V$ um espaço vetorial de dimensão $n$ sobre um corpo $K$. Seja $T$ um operador linear de $V$.
Prove que $T$ possui $n$ autovalores distintos se e somente se $T$ é diagonalizável e tem um vetor cíclico.

\subsection*{Demonstração}

Esta é uma proposição de equivalência, que exige a prova de duas implicações separadas.

\subsubsection*{Parte 1: Implicação Direta ($\Rightarrow$)}

\textbf{Hipótese:} O operador $T$ possui $n$ autovalores distintos em $K$, digamos
$\lambda_1, \lambda_2, \ldots, \lambda_n$. \\

\textbf{Tese:} O operador $T$ é diagonalizável e possui um vetor cíclico.

Vamos provar cada parte da tese separadamente.

\paragraph{a) Prova de que $T$ é diagonalizável.}
\begin{itemize}
    \item Sejam $v_1, v_2, \ldots, v_n$ autovetores não nulos correspondentes aos autovalores distintos
    $\lambda_1, \lambda_2, \ldots, \lambda_n$. Assim, $T(v_i) = \lambda_i v_i$ para cada $i$.
    \item Invocamos um teorema fundamental da teoria espectral: \textbf{Autovetores associados a
    autovalores distintos são linearmente independentes}.
    \item Portanto, o conjunto de autovetores $\{v_1, v_2, \ldots, v_n\}$ é um conjunto linearmente
    independente.
    \item Temos um conjunto de $n$ vetores linearmente independentes em um espaço vetorial $V$ de
    dimensão $n$. Por definição, este conjunto constitui uma \textbf{base} para $V$.
    \item Um operador linear é dito diagonalizável se, e somente se, existe uma base do espaço vetorial
    consistindo inteiramente de seus autovetores.
    \item Como acabamos de encontrar tal base, concluímos que \textbf{$T$ é diagonalizável}.
\end{itemize}

\paragraph{b) Prova de que $T$ possui um vetor cíclico.}
Um vetor $v \in V$ é dito cíclico para $T$ se o conjunto
$\{v, T(v), T^2(v), \ldots, T^{n-1}(v)\}$ forma uma base para $V$. Vamos construir explicitamente tal vetor.

\begin{itemize}
    \item Sejam $v_1, v_2, \ldots, v_n$ os autovetores linearmente independentes da parte (a).
    \item Considere o vetor $v$ definido pela soma destes autovetores:
    \[
        v = v_1 + v_2 + \cdots + v_n
    \]
    \item Vamos mostrar que o conjunto
    $\mathcal{B} = \{v, T(v), \ldots, T^{n-1}(v)\}$ é linearmente independente.
    \item Calculemos a ação de potências de $T$ sobre $v$:
    \[
        T^k(v) = T^k(v_1 + \cdots + v_n) = T^k(v_1) + \cdots + T^k(v_n) = \lambda_1^k v_1 + \cdots + \lambda_n^k v_n
    \]
    \item Agora, considere uma combinação linear nula dos vetores em $\mathcal{B}$:
    \[
        c_0 v + c_1 T(v) + \cdots + c_{n-1} T^{n-1}(v) = 0
    \]
    \item Substituindo as expressões para $T^k(v)$:
    \[
        \sum_{j=0}^{n-1} c_j T^j(v) =
        \sum_{j=0}^{n-1} c_j \left( \sum_{i=1}^n \lambda_i^j v_i \right) = 0
    \]
    \item Trocando a ordem dos somatórios (uma operação válida em somas finitas):
    \[
        \sum_{i=1}^n \left( \sum_{j=0}^{n-1} c_j \lambda_i^j \right) v_i = 0
    \]
    \item Seja $q(x) = c_0 + c_1 x + \cdots + c_{n-1}x^{n-1}$ um polinômio em $K[x]$. A equação acima pode ser
    escrita como:
    \[
        q(\lambda_1)v_1 + q(\lambda_2)v_2 + \cdots + q(\lambda_n)v_n = 0
    \]
    \item Como os autovetores $v_1, \ldots, v_n$ formam uma base, eles são linearmente independentes.
    Portanto, todos os coeficientes escalares nesta combinação linear devem ser nulos:
    \[
        q(\lambda_i) = 0 \quad \text{para todo } i = 1, \ldots, n
    \]
    \item Isto nos diz que o polinômio $q(x)$, que tem grau no máximo $n-1$, possui $n$ raízes distintas
    ($\lambda_1, \ldots, \lambda_n$).
    \item Um polinômio não nulo de grau $d$ pode ter no máximo $d$ raízes. A única forma de um polinômio
    de grau $\leq n-1$ ter $n$ raízes distintas é se ele for o polinômio nulo.
    \item Portanto, $q(x) = 0$, o que implica que todos os seus coeficientes são nulos:
    $c_0 = c_1 = \cdots = c_{n-1} = 0$.
    \item Isso prova que o conjunto $\mathcal{B}$ é linearmente independente. Sendo um conjunto com $n$
    vetores L.I. em um espaço de dimensão $n$, ele é uma base.
    \item Concluímos que \textbf{$v$ é um vetor cíclico para $T$}.
\end{itemize}

Assim, a implicação direta está provada.

\subsubsection*{Parte 2: Implicação Recíproca ($\Leftarrow$)}

\textbf{Hipótese:} O operador $T$ é diagonalizável e possui um vetor cíclico. \\

\textbf{Tese:} O operador $T$ possui $n$ autovalores distintos.

Procederemos por contradição.
\begin{itemize}
    \item Suponhamos que $T$ não possui $n$ autovalores distintos. Sejam $\mu_1, \mu_2, \ldots, \mu_d$
    os autovalores distintos de $T$, onde $d < n$.
    \item Analisemos as consequências de cada parte da hipótese:
    \begin{enumerate}
        \item \textbf{$T$ é diagonalizável:} Um operador é diagonalizável se, e somente se, seu polinômio
        minimal $m_T(x)$ se fatora em fatores lineares distintos. As raízes do polinômio minimal são
        precisamente os autovalores distintos do operador. Portanto:
        \[
            m_T(x) = (x - \mu_1)(x - \mu_2)\cdots(x - \mu_d)
        \]
        O grau do polinômio minimal é, consequentemente, $\deg(m_T) = d$.
        \item \textbf{$T$ possui um vetor cíclico:} Um teorema fundamental afirma que um operador possui
        um vetor cíclico se, e somente se, seu polinômio minimal é igual ao seu polinômio
        característico.
    \end{enumerate}
    \item Agora, confrontemos estas duas conclusões.
    \begin{itemize}
        \item Da primeira, temos que $\deg(m_T) = d$.
        \item Da segunda, temos que $m_T(x) = p_T(x)$. Sabemos que o grau do polinômio característico
        $p_T(x)$ é sempre igual à dimensão do espaço, que é $n$. Logo, $\deg(m_T) = \deg(p_T) = n$.
    \end{itemize}
    \item Isso nos leva a uma igualdade direta: $d = n$.
    \item No entanto, nossa suposição inicial para a prova por contradição foi que o número de
    autovalores distintos era $d < n$. A conclusão $d = n$ contradiz frontalmente a suposição
    $d < n$.
    \item A contradição surgiu da suposição de que $T$ não possuía $n$ autovalores distintos. Portanto,
    essa suposição deve ser falsa.
    \item Concluímos que \textbf{$T$ deve possuir $n$ autovalores distintos}.
\end{itemize}

A implicação recíproca está provada.

\subsection*{Conclusão}

Tendo provado ambas as implicações, estabelecemos a equivalência. Um operador sobre um espaço
de dimensão $n$ ter um espectro "completo" de $n$ autovalores distintos é a condição algébrica precisa
que garante simultaneamente a mais simples estrutura geométrica (ser diagonalizável) e a mais rica
estrutura dinâmica (possuir um vetor cíclico que gera todo o espaço sob a ação do operador).